\documentclass[10pt,a4paper]{amsart}
\usepackage[margin=19mm]{geometry}
\usepackage{amsmath,amssymb,mathtools}
\usepackage[scr=rsfs,cal=euler]{mathalfa}
\usepackage{xfrac}
\usepackage[unicode,hidelinks,bookmarks=false]{hyperref}
\newcommand{\ie}{i.e.\ }
\newcommand{\cf}{cf.\ }
\newcommand{\resp}{respectively\ }
\newcommand{\Subsection}{\subsection}
\providecommand{\nobreakdash}{\nobreak\hbox{-}}
\setlength{\emergencystretch}{2em}
\multlinegap0pt
\usepackage{mathtools}
\usepackage{amssymb}
\newtheorem{theorem}{Theorem}
\newcommand{\C}{\mathbb{C}}
\newcommand{\Der}{\mathrm{Der}}
\newcommand{\Z}{\mathbb{Z}}
\newcommand{\hDer}{\widehat{\Der(A)}}
\newcommand{\ol}[1]{\overline{#1}}
\newcommand{\pa}{\partial}
\title{Orthogonal polynomials and the centres of the universal central extensions of superelliptic Krichever--Novikov algebras}
\author{Felipe Albino dos Santos}
\date{Source: arXiv:2605.02530v2 (CC BY 4.0)}
\begin{document}
\maketitle

\noindent\emph{Source and licence.} Orthogonal polynomials and the centres of the universal central extensions of superelliptic Krichever--Novikov algebras. Version arXiv:2605.02530v2 (CC BY 4.0), distributed by arXiv under the Creative Commons Attribution 4.0 International licence, http://creativecommons.org/licenses/by/4.0/, as recorded in the arXiv licence metadata for this exact version and shown on the abstract page https://arxiv.org/abs/2605.02530v2. Full licence notice: \texttt{licenses/arxiv-2605.02530-CC-BY-4.0.txt}.\par\smallskip

\noindent\textit{Editorial scope.} Counted unit: the article's theorem thm:cross\_closed (source0.tex lines 905--920). The article's complete original proof of it is delivered with it (source0.tex lines 922--955, one \texttt{proof} environment of 34 lines). No mathematics is written, repaired or re-derived: the statement and the proof are the article's own lines, in the article's own notation, and no retained line is substituted or re-flowed.  Retained uncounted as the source-local context the proof needs: lines 450--452; lines 599--602; lines 643--645; lines 694--714.  Nothing was omitted from any retained span, so the package carries no omission record.  No retained line contains a citation, so the wrapper carries no bibliography and no bibliography entry.  No retained line contains a figure environment, an image-inclusion command or drawing code, so no image is delivered and the package declares no asset dependency.  Editorial formatting only, in addition: an AMS \texttt{amsart} preamble that restates the article's own declarations and macros used by the retained lines, namely the environments \texttt{theorem} on separate counters, together with the standard packages those lines need.  The retained environments print in this document as Theorem 1 in order of appearance, whereas the article prints the same items with the numbering of its own full document.  The complete native source of the article as distributed in the arXiv e-print for this exact version is bundled unchanged as \texttt{sources/199720/source0.tex}.
\begin{equation}\label{eq:psi_def}
  \psi(f\pa,\,g\pa) = \ol{\pa(f)\,\pa^2(g)} .
\end{equation}

\begin{equation}\label{eq:legendre_red}
  [x^k] = p_k(a)\,\omega_2 \quad (k\geq 0),\qquad
  [x^{-k}] = p_{k-1}(a)\,\omega_1 \quad (k\geq 1),
\end{equation}

\begin{equation}\label{eq:cross_general}
  \psi(e_k,f_s) = k\!\sum_{j=-3}^{1} C_j(s,a)\,[x^{n+j}]_{A/\pa A},
\end{equation}

\begin{theorem}\label{thm:antiderivative}
Let $m=2$ and $P(x)=x^2-2ax+1$ with $a\neq\pm1$, let $\pa=u\,d/dx$ and let
$\psi$ be the cocycle \eqref{eq:psi_def} of $\hDer$, with $e_k=x^k\pa$,
$f_s=x^su\pa$ and $\omega_2=[1]$ as above. Let $s\in\Z$ and put
$n=s+1\geq1$. Then
\begin{equation}\label{eq:antideriv}
  \psi(e_1,f_s) = \frac{p_{n-1}(a)-p_{n+1}(a)}{2n+1}\,\omega_2.
\end{equation}
Moreover:
\begin{equation}\label{eq:antideriv_integral}
  \frac{p_{n-1}(a)-p_{n+1}(a)}{2n+1} = \int_a^1 p_n(t)\,dt.
\end{equation}
Thus the cross-cocycle $\psi(e_1,f_s)$ equals the definite integral
of the $n$-th Legendre polynomial from $a$ to $1$, with
$n=s+1$. We write
\begin{equation}\label{eq:In}
  I_n(a)\;\coloneqq\;\int_a^1p_n(t)\,dt\;=\;\frac{p_{n-1}(a)-p_{n+1}(a)}{2n+1},
  \qquad n\geq1 ,
\end{equation}
for this scalar.
\end{theorem}

\begin{theorem}\label{thm:cross_closed}
Let $m=2$ and $P(x)=x^2-2ax+1$ with $a\neq\pm1$, and let $\psi$, $e_k$, $f_s$,
$\omega_2$ be as in Theorem~\ref{thm:antiderivative}. Let $k,s\in\Z$ with
$k\neq0$ and put $n=k+s$. If $n\geq3$ then
\begin{equation}\label{eq:cross_closed}
  \psi(e_k,f_s)
  \;=\;
  \frac{k}{2n-1}\Bigl(\bigl((3k-1)s+2k-1\bigr)I_n(a)
      \;-\;3(k-1)s\,I_{n-2}(a)\Bigr)\,\omega_2 .
\end{equation}
In particular the value lies on the line $\C\omega_2$ for every index once
$n\geq3$, it vanishes at $a=1$, and the second term is absent exactly when
$k=1$ or $s=0$. For $k=1$ the formula reduces to
$\psi(e_1,f_s)=I_n(a)\,\omega_2$, which is the $n\geq3$ case of
Theorem~\ref{thm:antiderivative}.
\end{theorem}

\begin{proof}
The five exponents occurring in \eqref{eq:cross_general} are $n-3,\dots,n+1$,
and $n\geq3$ makes all of them non-negative, so the Legendre reduction
\eqref{eq:legendre_red} applies to each and gives
\[
  \psi(e_k,f_s)=k\Bigl(\sum_{j=-3}^{1}C_j(s,a)p_{n+j}(a)\Bigr)\omega_2 .
\]
The coefficients $C_j$ are polynomial of degree at most $2$ in $a$, so the sum
is brought back to the Legendre basis by applying
$ap_\mu=\bigl((\mu+1)p_{\mu+1}+\mu p_{\mu-1}\bigr)/(2\mu+1)$ at most twice to
each term. Collecting, the coefficients of $p_n$ and of $p_{n-2}$ cancel
identically and what survives is
\[
  \psi(e_k,f_s)=\bigl(\alpha\,p_{n+1}+\beta\,p_{n-1}+\delta\,p_{n-3}\bigr)\omega_2 ,
\]
with
\[
  \alpha=\frac{k\bigl(3k^2-3k(n+1)+n+1\bigr)}{4n^2-1},\qquad
  \delta=\frac{3k(k-1)(k-n)}{4n^2-8n+3},\qquad
  \beta=-\alpha-\delta .
\]
The identity $\alpha+\beta+\delta=0$ is what lets the three-term combination
collapse onto two antiderivatives; it is an identity of rational functions in
$k$ and $n$, obtained by direct expansion. Using it,
\[
  \psi(e_k,f_s)
   =\bigl(\alpha(p_{n+1}-p_{n-1})+\delta(p_{n-3}-p_{n-1})\bigr)\omega_2
   =\bigl(-(2n+1)\alpha\,I_n+(2n-3)\delta\,I_{n-2}\bigr)\omega_2 ,
\]
and substituting $s=n-k$ into the displayed $\alpha$ and $\delta$ turns
$-(2n+1)\alpha$ into $k\bigl((3k-1)s+2k-1\bigr)/(2n-1)$ and $(2n-3)\delta$ into
$-3k(k-1)s/(2n-1)$, which is \eqref{eq:cross_closed}. Setting $k=1$ kills
$\delta$ and leaves $-(2n+1)\alpha=1$.
\end{proof}

\end{document}
